\subsection{The convexity theorem}

THEOREM 1. --- Let X be a locally compact space, \(\mu\) a positive measure on X, F a real Banach space, D a closed convex set in F, and f a function on X such that \(\mathbf{f}(\mathbf{X}) \subset \mathbf{D}\) . For every non-negligible integrable numerical function \(g \geqslant 0\) such that \(\mathbf{fg}\) is integrable, the point \(\frac{\int \mathbf{fg} d\mu}{\int g d\mu}\) belongs to D.

For, let \(\mathbf{F}'\) be the dual of \(\mathbf{F}\) and let \(\langle \mathbf{z},\mathbf{a}'\rangle \leqslant \alpha (\mathbf{a}'\in \mathbf{F}',\alpha \in \mathbf{R})\) be a relation defining a closed half-space containing D. Since \(\mathbf{fg}\) is integrable, so is the numerical function \(\langle \mathbf{fg},\mathbf{a}'\rangle = \langle \mathbf{f},\mathbf{a}'\rangle g\) , and

\[
\int \left\langle \mathbf {f} g, \mathbf {a} ^ {\prime} \right\rangle d \mu = \left\langle \int \mathbf {f} g d \mu , \mathbf {a} ^ {\prime} \right\rangle
\]

(§4, No.~2, Cor. 1 of Th. 1); but, by hypothesis, \(\langle \mathbf{f}(x), \mathbf{a}' \rangle \leqslant \alpha\) for all \(x \in X\) , therefore \(\langle \mathbf{f}(x)g(x), \mathbf{a}' \rangle \leqslant \alpha g(x)\) ; on integrating, we have

\[
\left\langle \int \mathbf {f} g d \mu , \mathbf {a} ^ {\prime} \right\rangle \leqslant \alpha \int g d \mu .
\]

This proves that the point \(\frac{\int\mathbf{f}g d\mu}{\int g d\mu}\) belongs to every closed half-space containing D; but, by the Hahn--Banach theorem, D is the intersection of the closed half-spaces containing it (TVS, II, §5, No.~3, Cor. 1 of Prop. 4), whence the theorem.

COROLLARY. --- If the positive measure \(\mu\) has total mass equal to 1 and if \(\mathbf{f}\) is integrable, then \(\int \mathbf{f} d\mu\) belongs to the closed convex envelope of \(\mathbf{f}(\mathbf{X})\) in \(\mathbf{F}\) .

It suffices to take for g the constant function 1.

